\documentclass[11pt]{article}
\usepackage[margin=1in]{geometry}
\usepackage{amsmath,amssymb,amsthm,booktabs,graphicx,hyperref,microtype}
\hypersetup{colorlinks=true,linkcolor=blue,citecolor=blue,urlcolor=blue}

\newtheorem{theorem}{Theorem}
\newtheorem{lemma}{Lemma}
\newcommand{\E}{\mathcal E}
\newcommand{\N}{\mathbb N}

\title{An exact nonlinear dual formula on the two-site birth--death chain\\
\large A substantial partial answer to Open Question 4 of arXiv:math/0101227}
\author{Autonomous mathematical discovery run \texttt{fa\_banach\_001}, lane 8}
\date{12 August 2026}

\begin{document}
\maketitle

\noindent\textbf{Status.}
Candidate substantial partial result, likely valid.  For every $p>2$ we
prove a complete nonlinear sup--inf/inf--sup formula for the sharp
$\ell^p$ Sobolev constant of the first nontrivial Dirichlet birth--death
chain.  This is the Sobolev partner of a Nash inequality, but it is not the
centered ergodic Nash or logarithmic-Sobolev constant asked for in full by
the source.

\section{Source question and later-literature boundary}
Chen asks for a one-dimensional Nash criterion and, separately, for a
variational formula in the completed dual form of his Theorem 4.1(2) for a
Nash or logarithmic-Sobolev inequality \cite[p.~6, Open Problem 3 and Open
Question 4]{Chen01}.  The exact source passage is reproduced below.

\begin{figure}[ht]
  \centering
  \includegraphics[width=\textwidth]{figures/open_questions_crop.png}
  \caption{Open Problem 3 and Open Question 4 in the source PDF, page 6.}
\end{figure}

Mao subsequently gave if-and-only-if Nash criteria for birth--death
processes and diffusions on the line, explicitly identifying Chen's question
\cite{Mao02}.  Thus Problem 3 is already answered.  Chen later extended the
Poincar\'e formulas to Banach/Orlicz norms \cite{Chen02,Chen03}.  Those results
give exact criteria, explicit constant-factor bounds, and inf--sup formulas
for \emph{upper bounds}; they do not state the exact dual equality requested
for the centered Nash or logarithmic-Sobolev constant.  The result below
isolates a setting in which the exact nonlinear equality can be recovered.

\section{Proof intuition}
The derivative of an $L^p$ Sobolev quotient is nonlinear: its source is
$u^{p-1}$, not $u$.  On a one-dimensional Dirichlet chain, summing the
Euler--Lagrange equation from a site to the Neumann endpoint converts that
source into a positive tail.  Dividing the tail by the edge flux gives a
nonlinear analogue of Chen's $I$-functional.  A summation-by-parts identity
proves the sup--inf half on every finite chain.  For two sites, the two tail
functionals are respectively strictly decreasing and strictly increasing in
the only projective coordinate, so their crossing proves the dual inf--sup
half as well.

\section{A general exact sup--inf half-formula}
Let $n\ge1$, let masses $\mu_i>0$ and conductances $c_i>0$, and put $f_0=0$.
For $p>2$, define
\[
 \|f\|_{p,\mu}=\left(\sum_{i=1}^n\mu_i|f_i|^p\right)^{1/p},\qquad
 \E(f)=\sum_{i=1}^n c_i(f_i-f_{i-1})^2,
\]
and the sharp Dirichlet Sobolev constant
\[
 A_p=\max_{f\ne0}\frac{\|f\|_{p,\mu}^2}{\E(f)}.
\]
For $f_0=0<f_1<\cdots<f_n$, set
\begin{equation}\label{eq:Jgeneral}
 J_i(f)=
 \frac{\sum_{j=i}^n\mu_j f_j^{p-1}}
 {c_i(f_i-f_{i-1})\,\|f\|_{p,\mu}^{p-2}}.
\end{equation}

\begin{lemma}\label{lem:half}
For every finite Dirichlet birth--death chain and every $p>2$,
\[
 A_p=\sup_{0=f_0<f_1<\cdots<f_n}\ \min_{1\le i\le n}J_i(f).
\]
\end{lemma}

\begin{proof}
Let $m=\min_iJ_i(f)$.  Multiply the defining inequality for $m$ by
$f_i-f_{i-1}$ and sum over $i$.  Reversing the finite sums gives
\begin{align*}
 m\|f\|_{p,\mu}^{p-2}\E(f)
 &\le \sum_{i=1}^n(f_i-f_{i-1})
       \sum_{j=i}^n\mu_j f_j^{p-1}\\
 &=\sum_{j=1}^n\mu_jf_j^{p-1}
       \sum_{i=1}^j(f_i-f_{i-1})
 =\|f\|_{p,\mu}^{p}.
\end{align*}
Hence $m\le\|f\|_{p,\mu}^2/\E(f)\le A_p$.

The maximum defining $A_p$ exists.  Replacing an optimizer by its absolute
value cannot increase the energy.  Its Euler equation and the positive
inverse of the Dirichlet stiffness matrix then make it strictly positive;
summing the Euler equation from $i$ to $n$ makes it strictly increasing.
For such an optimizer $u$,
\[
 c_i(u_i-u_{i-1})
 =\frac{\E(u)}{\|u\|_{p,\mu}^p}
   \sum_{j=i}^n\mu_j u_j^{p-1}.
\]
Thus every $J_i(u)=\|u\|_{p,\mu}^2/\E(u)=A_p$, proving the reverse
inequality and the lemma.
\end{proof}

\section{The complete two-site dual formula}
Take $n=2$, $\mu_1=\mu_2=c_1=c_2=1$.  Thus
\[
 \E(x)=x_1^2+(x_2-x_1)^2,\qquad
 A_p=\max_{x\ne0}\frac{(|x_1|^p+|x_2|^p)^{2/p}}
 {x_1^2+(x_2-x_1)^2}.
\]
For $0<s<1$, put
\begin{align}
 Q_p(s)&=\frac{(1+s^p)^{2/p}}{s^2+(1-s)^2},\label{eq:Q}\\
 J_1(s)&=\frac{1+s^{p-1}}{s(1+s^p)^{(p-2)/p}},\qquad
 J_2(s)=\frac{1}{(1-s)(1+s^p)^{(p-2)/p}}.\label{eq:Jtwo}
\end{align}

\begin{theorem}\label{thm:main}
For every $p>2$, there is a unique $s_p\in(0,1)$ satisfying
\begin{equation}\label{eq:root}
 1-2s_p+s_p^{p-1}-s_p^p=0.
\end{equation}
The sharp two-site constant has the completed nonlinear dual formula
\begin{equation}\label{eq:dual}
 \boxed{\quad
 A_p=\sup_{0<s<1}\min\{J_1(s),J_2(s)\}
 =\inf_{0<s<1}\max\{J_1(s),J_2(s)\}
 =Q_p(s_p)=J_1(s_p)=J_2(s_p).
 \quad}
\end{equation}
\end{theorem}

\begin{proof}
Let
\[
 H_p(s)=1-(p-2)s^{p-1}+(p-1)s^p.
\]
Its minimum on $[0,1]$ occurs at $r=(p-2)/p$ and equals
$1-r^p>0$.  Direct logarithmic differentiation of \eqref{eq:Jtwo} gives
\[
 J_1'(s)<0,\qquad J_2'(s)>0,
\]
with the two derivative numerators equal to $-H_p(s)$ and $H_p(s)$,
respectively, after multiplication by positive denominators.  Moreover,
$J_1(s)=J_2(s)$ is exactly equation \eqref{eq:root}.  Since $J_1$ starts
above $J_2$ at $0$ and ends below it at $1$, there is one crossing $s_p$.
The minimum of the maximum and the maximum of the minimum therefore both
occur at this crossing.

Replacing $(x_1,x_2)$ by $(|x_1|,|x_2|)$ cannot increase the denominator of
the Sobolev quotient.  After scaling $x_2=1$, its relevant projective
coordinate is $s=x_1/x_2\ge0$.  On $(0,\infty)$,
\[
 \frac{Q_p'(s)}{Q_p(s)}=
 \frac{2(1-2s+s^{p-1}-s^p)}{(1+s^p)(s^2+(1-s)^2)}.
\]
The sign change just proved shows that $Q_p$ increases up to $s_p$ and
decreases thereafter (for $s\ge1$ the displayed numerator is negative).
Thus $A_p=Q_p(s_p)$.  Finally, equation \eqref{eq:root} implies
\[
 s_p^2+(1-s_p)^2=(1-s_p)(1+s_p^p),
\]
so $Q_p(s_p)=J_2(s_p)=J_1(s_p)$.  This proves \eqref{eq:dual}.
\end{proof}

\section{The simpler Banach tail formula is strictly nonsharp}
The Banach-space upper formula in Chen's later theory replaces the nonlinear
tail $f^{p-1}$ by the norm of an ordinary tail.  At $p=4$, the relevant
Banach norm on $f^2$ is $\ell^2$.  Its direct two-site tail functional is
\begin{equation}\label{eq:U}
 U_4=\inf_{t>1}\max\left\{\sqrt{1+t^2},\frac{t}{t-1}\right\}.
\end{equation}

\begin{theorem}\label{thm:gap}
For the two-site chain,
\[
 A_4<2.08<2.13<U_4.
\]
In particular, the upper formula \eqref{eq:U} cannot be upgraded to an exact
identity even in this smallest nontrivial model.
\end{theorem}

\begin{proof}
The branches in \eqref{eq:U} are strictly increasing and decreasing,
respectively.  They meet at the unique root $t_*$ of
\[
 t^4-2t^3+t^2-2t+1=0.
\]
Exact rational evaluation gives $1.883<t_*<1.884$, hence
\[
 U_4^2>1+(1.883)^2=4.545689>(2.13)^2.
\]

Equation \eqref{eq:root} at $p=4$ is equivalent to
$s^4-s^3+2s-1=0$.  Exact evaluation gives
$0.535<s_4<0.536$.  On this interval the numerator $1+s^4$ and denominator
$2s^2-2s+1$ are increasing, so
\[
 A_4^2<
 \frac{1+(0.536)^4}{(2(0.535)^2-2(0.535)+1)^2}
 =\frac{270634747904}{63114000625}
 <(2.08)^2.
\]
All signs and rational comparisons are independently reproduced by the
supplied verifier.
\end{proof}

\section{Verification, novelty bounds, and limitations}
Run \texttt{code/verify\_two\_site.py} from the repository root.  It checks
the $p=4$ derivative factorizations symbolically, certifies both rational root
brackets and the strict $0.05$ separation, and evaluates the general formula
at five exponents with 60-decimal arithmetic.  The proof for arbitrary
$p>2$ is algebraic and does not depend on those numerical checks.

The bounded novelty search covered the run registry and attempt indexes,
Chen's 2001 source, Mao's 2002 Nash paper, Chen's 2002/2003 Banach-space and
birth--death papers, the later monograph chapter summarizing those results,
and web/arXiv searches combining ``Collatz--Wielandt,'' ``Nash,'' ``Sobolev
constant,'' ``birth--death,'' and ``variational formula.''  General nonlinear
Perron--Frobenius theory is known, but no source stating formulas
\eqref{eq:dual} or the strict comparison in Theorem~\ref{thm:gap} was found.
This is bounded evidence only; the specialization may be folklore.

The result is deliberately scoped.  The centered Neumann Nash quotient has
an Euler source containing both $f$ and a constant from its $L^1$ factor;
the logarithmic-Sobolev source contains
$f\log(f^2/\|f\|_2^2)$.  Both change sign, destroying the positive-tail
ordering used here.  A reviewer should verify that this distinction is
maintained and that the packet is not read as a solution of the full source
question.

\small
\begin{thebibliography}{9}
\bibitem{Chen01}
M.-F. Chen,
\emph{Explicit criteria for several types of ergodicity},
Chinese J. Appl. Probab. Statist. 17 (2001), 1--8;
\href{https://arxiv.org/abs/math/0101227}{arXiv:math/0101227}.

\bibitem{Mao02}
Y.-H. Mao,
\emph{Nash inequalities for Markov processes in dimension one},
Acta Math. Sin. (Engl. Ser.) 18 (2002), 147--156,
\href{https://doi.org/10.1007/s101140100128}{doi:10.1007/s101140100128}.

\bibitem{Chen02}
M.-F. Chen,
\emph{Variational formulas of Poincar\'e-type inequalities in Banach spaces
of functions on the line},
Acta Math. Sin. (Engl. Ser.) 18 (2002), 417--436,
\href{https://doi.org/10.1007/s10114-002-0176-8}{doi:10.1007/s10114-002-0176-8}.

\bibitem{Chen03}
M.-F. Chen,
\emph{Variational formulas of Poincar\'e-type inequalities for birth--death
processes},
Acta Math. Sin. (Engl. Ser.) 19 (2003), 625--644,
\href{https://doi.org/10.1007/s10114-003-0308-9}{doi:10.1007/s10114-003-0308-9}.
\end{thebibliography}
\normalsize

\end{document}
